% GENERATED FILE. Edit canonical lessons, then run npm run generate:books.
% canonical-source-hash: fnv1a64:7fdc6af58bfa5d01
% canonical-lessons: LA-C337-palma, LA-C337-praemium, LA-C337-certo, LA-C337-certamen, LA-C337-plausus

\chapter{Palma et praemium}
\label{ch:la-palma-et-praemium}
\hlchaptermodality{337}

\begin{chapteropening}
\textbf{By the end of this chapter:} \emph{I can tell who won a contest: the palm of victory, the prize, to compete, the contest and the applause.}

Complete the last lesson of chapter 337: I can tell who won a contest: the palm of victory, the prize, to compete, the contest and the applause.
\end{chapteropening}

\section[palma, palmae]{palma, palmae — the palm of victory; the palm of the hand; a palm tree}

\label{lesson:LA-C337-palma}



\begin{quote}
Before the new one: say the Latin for a hunt, then the Latin for a beast.
\end{quote}

\subsection*{You'll want to know: palma, palmae}
\textbf{palma, palmae} — \textquotedblleft{}the palm of victory; the palm of the hand; a palm tree\textquotedblright{}.

\begin{quote}
Aurīga palmam et corōnam accēpit.
\end{quote}

\begin{quote}
\emph{The charioteer received the palm and a crown.}
\end{quote}

Feminine. The flat of the hand, then the palm tree, whose branch was given to winners.

\begin{cousinweb}
\textbf{palma} is inherited from a Proto-Indo-European form for \textquotedblleft{}the flat of the hand\textquotedblright{} (EtymDB), matching Greek \emph{palamē}. The palm tree was named from its hand-like leaves. English \emph{palm} comes from it.
\end{cousinweb}

\begin{grammarlens}[title={What you've built}]
One more word for telling, explaining and describing at B1.
\end{grammarlens}

\subsection*{Your turn}
\begin{itemize}
\raggedright
  \item \emph{Say it:} \emph{palma, palmae}
  \item \emph{Say it:} \emph{palma, palmae}, once more
  \item \emph{Say it:} a sentence of your own with \emph{palma, palmae}
  \item \emph{Recall:} say the Latin for a hunt, then the Latin for a beast, then say \emph{palma, palmae} again
\end{itemize}

\begin{tcolorbox}[breakable,colback=teal!4,colframe=teal!35!black,title={Before you move on}]
What does palma, palmae mean? (\textquotedblleft{}The palm of victory; the palm of the hand; a palm tree\textquotedblright{}.) Say it once more.
\end{tcolorbox}
\section[praemium, praemiī]{praemium, praemiī — a prize, a reward}

\label{lesson:LA-C337-praemium}



\begin{quote}
Before the new one: say the Latin for a beast, then the Latin for the palm of victory.
\end{quote}

\subsection*{You'll want to know: praemium, praemiī}
\textbf{praemium, praemiī} — \textquotedblleft{}a prize, a reward\textquotedblright{}.

\begin{quote}
Praemium magnum erat: centum dēnāriī.
\end{quote}

\begin{quote}
\emph{The prize was large: a hundred denarii.}
\end{quote}

Neuter. First booty taken before others, then any reward. English \emph{premium}.

\begin{cousinweb}
\textbf{praemium} is built on \emph{prae}, \textquotedblleft{}before\textquotedblright{}, and \emph{emō}, \textquotedblleft{}to take, to buy\textquotedblright{} (EtymDB, Lewis \& Short): what one takes before others. English \emph{premium} comes from it.
\end{cousinweb}

\begin{grammarlens}[title={What you've built}]
One more word for telling, explaining and describing at B1.
\end{grammarlens}

\subsection*{Your turn}
\begin{itemize}
\raggedright
  \item \emph{Say it:} \emph{praemium, praemiī}
  \item \emph{Say it:} \emph{praemium, praemiī}, once more
  \item \emph{Say it:} a sentence of your own with \emph{praemium, praemiī}
  \item \emph{Recall:} say the Latin for a beast, then the Latin for the palm of victory, then say \emph{praemium, praemiī} again
\end{itemize}

\begin{tcolorbox}[breakable,colback=teal!4,colframe=teal!35!black,title={Before you move on}]
What does praemium, praemiī mean? (\textquotedblleft{}A prize, a reward\textquotedblright{}.) Say it once more.
\end{tcolorbox}
\section[certō, certāre]{certō, certāre — to compete, to contend}

\label{lesson:LA-C337-certo}



\begin{quote}
Before the new one: say the Latin for the palm of victory, then the Latin for a prize.
\end{quote}

\subsection*{You'll want to know: certō, certāre}
\textbf{certō, certāre} — \textquotedblleft{}to compete, to contend\textquotedblright{}.

\begin{quote}
Duo aurīgae dē palmā certābant.
\end{quote}

\begin{quote}
\emph{Two charioteers were competing for the palm.}
\end{quote}

\emph{certō, certāre, certāvī, certātum}. Do not confuse the adverb \emph{certō}, 'for certain'.

\begin{cousinweb}
\textbf{certō} is built on \emph{cernō}, \textquotedblleft{}to decide\textquotedblright{} (Lewis \& Short): to decide something by a contest.
\end{cousinweb}

\begin{grammarlens}[title={What you've built}]
One more word for telling, explaining and describing at B1.
\end{grammarlens}

\subsection*{Your turn}
\begin{itemize}
\raggedright
  \item \emph{Say it:} \emph{certō, certāre}
  \item \emph{Say it:} \emph{certō, certāre}, once more
  \item \emph{Say it:} a sentence of your own with \emph{certō, certāre}
  \item \emph{Recall:} say the Latin for the palm of victory, then the Latin for a prize, then say \emph{certō, certāre} again
\end{itemize}

\begin{tcolorbox}[breakable,colback=teal!4,colframe=teal!35!black,title={Before you move on}]
What does certō, certāre mean? (\textquotedblleft{}To compete, to contend\textquotedblright{}.) Say it once more.
\end{tcolorbox}
\section[certāmen, certāminis]{certāmen, certāminis — a contest, a struggle}

\label{lesson:LA-C337-certamen}



\begin{quote}
Before the new one: say the Latin for a prize, then the Latin for to compete.
\end{quote}

\subsection*{You'll want to know: certāmen, certāminis}
\textbf{certāmen, certāminis} — \textquotedblleft{}a contest, a struggle\textquotedblright{}.

\begin{quote}
Certāmen longum et difficile erat.
\end{quote}

\begin{quote}
\emph{The contest was long and hard.}
\end{quote}

Neuter, 3rd declension, from \emph{certō}. Any contest, friendly or hostile.

\begin{cousinweb}
\textbf{certāmen} is built on \emph{certō}, \textquotedblleft{}to contend\textquotedblright{} (Lewis \& Short), with the ending \emph{-men}.
\end{cousinweb}

\begin{grammarlens}[title={What you've built}]
One more word for telling, explaining and describing at B1.
\end{grammarlens}

\subsection*{Your turn}
\begin{itemize}
\raggedright
  \item \emph{Say it:} \emph{certāmen, certāminis}
  \item \emph{Say it:} \emph{certāmen, certāminis}, once more
  \item \emph{Say it:} a sentence of your own with \emph{certāmen, certāminis}
  \item \emph{Recall:} say the Latin for a prize, then the Latin for to compete, then say \emph{certāmen, certāminis} again
\end{itemize}

\begin{tcolorbox}[breakable,colback=teal!4,colframe=teal!35!black,title={Before you move on}]
What does certāmen, certāminis mean? (\textquotedblleft{}A contest, a struggle\textquotedblright{}.) Say it once more.
\end{tcolorbox}
\section[plausus, plausūs]{plausus, plausūs — applause, clapping}

\label{lesson:LA-C337-plausus}



\begin{quote}
Before the new one: say the Latin for to compete, then the Latin for a contest.
\end{quote}

\subsection*{You'll want to know: plausus, plausūs}
\textbf{plausus, plausūs} — \textquotedblleft{}applause, clapping\textquotedblright{}.

\begin{quote}
Magnus plausus in theātrō erat.
\end{quote}

\begin{quote}
\emph{There was loud applause in the theatre.}
\end{quote}

Masculine, 4th declension, from \emph{plaudō}, 'to clap' (taught). Plays ended with the call \emph{plaudite!}

\begin{cousinweb}
\textbf{plausus} is built on \emph{plaudō}, \textquotedblleft{}to clap\textquotedblright{} (Lewis \& Short). English \emph{applause} comes from its compound \emph{applaudō}.
\end{cousinweb}

\begin{grammarlens}[title={What you've built}]
One more word for telling, explaining and describing at B1.
\end{grammarlens}

\subsection*{Your turn}
\begin{itemize}
\raggedright
  \item \emph{Say it:} \emph{plausus, plausūs}
  \item \emph{Say it:} \emph{plausus, plausūs}, once more
  \item \emph{Say it:} a sentence of your own with \emph{plausus, plausūs}
  \item \emph{Recall:} say the Latin for to compete, then the Latin for a contest, then say \emph{plausus, plausūs} again
\end{itemize}

\begin{tcolorbox}[breakable,colback=teal!4,colframe=teal!35!black,title={Before you move on}]
What does plausus, plausūs mean? (\textquotedblleft{}Applause, clapping\textquotedblright{}.) Say it once more.
\end{tcolorbox}
